\documentclass[10pt]{article}

\usepackage[margin=1in]{geometry}
\usepackage{booktabs}
\usepackage{longtable}
\usepackage{amsmath}
\usepackage{graphicx}
\usepackage{svg}
\usepackage{float}
\usepackage[hidelinks]{hyperref}

\graphicspath{{../results/figures/}}
\svgpath{{../results/figures/}}

\title{Supplementary Material: Household and Community Heterogeneity in Childhood Stunting in Ghana}
\author{Gideon Ofosu Kyere\\\small Kwame Nkrumah University of Science and Technology, Kumasi-Ghana\\\small ORCID: 0009-0003-9848-8437}
\date{Preprint revision --- 9 October 2026}

% Supporting-information numbering
\renewcommand{\thetable}{S\arabic{table}}
\renewcommand{\thefigure}{S\arabic{figure}}

\begin{document}
\maketitle

\section{Analytic sample validation}

The reconstructed analytic sample contained 4,928 children, including 927 who were stunted, across 3,545 households and 616 survey communities. The rescaled DHS sampling weights summed to 4,292.97, reproducing the official weighted height-for-age denominator of 4,293 after rounding. Survey-weighted stunting prevalence was 17.39\%.

\section{Household structure}

Most households contributed a single eligible child: 2,398 of 3,545 households had one child in the analytic sample. A further 951 contributed two children, 165 contributed three, 26 contributed four, four contributed five, and one contributed nine. Overall, 1,147 households contained more than one eligible child, and about 51.3\% of analysed children lived in these multi-child households.

\section{Extended Model 3 diagnostics}

The strengthened Model 3 diagnostic run used eight independent chains, each with 500 warmup and 500 retained draws, for 4,000 posterior draws.

\begin{table}[H]
\centering
\caption{Sampler diagnostics for the strengthened Model 3 run}
\begin{tabular}{rrrrr}
\toprule
Chain & Divergences & BFMI & Max tree depth & Mean acceptance \\
\midrule
1 & 0 & 0.519 & 6 & 0.946 \\
2 & 0 & 0.655 & 6 & 0.943 \\
3 & 0 & 0.530 & 6 & 0.935 \\
4 & 0 & 0.660 & 6 & 0.935 \\
5 & 0 & 0.514 & 6 & 0.937 \\
6 & 0 & 0.562 & 6 & 0.937 \\
7 & 0 & 0.534 & 6 & 0.941 \\
8 & 0 & 0.651 & 6 & 0.929 \\
\bottomrule
\end{tabular}
\end{table}

No chain reached the configured maximum tree depth. Selected fixed-effect diagnostics were close to 1.00, but the locked community and household SDs had $\hat R$ approximately 1.021 and 1.018 and bulk ESS approximately 544 and 524. Both exceed the commonly used 1.01 guideline, so the locked fit does not meet the strict joint-posterior screen. The stored household-versus-community pattern remains exploratory; this is a limitation on validated inference, not only on exact interval endpoints. The separate reproduction audit below does not replace or certify this eight-chain posterior.

\section{Variance decomposition}

\begin{table}[H]
\centering
\caption{Retained exploratory variance summaries from the locked eight-chain Model 3}
\begin{tabular}{lrrr}
\toprule
Quantity & Median & 2.5th percentile & 97.5th percentile \\
\midrule
Community SD & 0.389 & 0.067 & 0.610 \\
Household SD & 1.234 & 0.940 & 1.535 \\
Community ICC & 0.030 & 0.001 & 0.073 \\
Household VPC & 0.307 & 0.203 & 0.406 \\
Same-household correlation & 0.339 & 0.236 & 0.436 \\
Community MOR & 1.45 & 1.07 & 1.79 \\
Household MOR & 3.24 & 2.45 & 4.32 \\
\bottomrule
\end{tabular}
\end{table}

\section{Prior sensitivity}

Historical selected fixed-effect summaries were similar under tighter and wider weakly regularizing priors, but both saved reproduction fits failed the strict full-parameter screen. These comparisons remain exploratory. The posterior median OR for unimproved water was about 1.48 under both alternatives, compared with 1.50 in the primary model. For higher maternal education, the corresponding ORs were about 0.44 under the tighter prior, 0.37 in the strengthened primary run, and 0.35 under the wider prior.


\begin{table}[H]
\centering
\caption{Prior-sensitivity summaries for selected fixed effects}
\small
\begin{tabular}{llrrr}
\toprule
Prior & Comparison & Median OR & 2.5th percentile & 97.5th percentile \\
\midrule
Tighter & Male vs female & 1.45 & 1.21 & 1.76 \\
Tighter & Unimproved water & 1.48 & 1.14 & 1.88 \\
Tighter & Unimproved sanitation & 1.15 & 0.91 & 1.47 \\
Tighter & Higher maternal education vs none & 0.44 & 0.26 & 0.70 \\
Wider & Male vs female & 1.52 & 1.25 & 1.85 \\
Wider & Unimproved water & 1.48 & 1.12 & 2.06 \\
Wider & Unimproved sanitation & 1.10 & 0.84 & 1.43 \\
Wider & Higher maternal education vs none & 0.35 & 0.18 & 0.66 \\
\bottomrule
\end{tabular}
\begin{minipage}{0.98\textwidth}\footnotesize \textit{Note.} Retained historical summaries; both available prior reproduction fits fail the strict full-parameter screen.\end{minipage}
\end{table}

\section{Survey-weight sensitivity}

For the survey-weight sensitivity analysis, DHS household-member weights were rescaled to have mean one and then used to weight the Bernoulli log-likelihood contributions. The water estimate became less precise (OR 1.38, 95\% posterior interval 0.90--2.07), although the posterior probability that its coefficient was positive remained about 0.94. The directions of the sex, age, richest-versus-poorest wealth, and higher maternal-education estimates were unchanged.


\begin{table}[H]
\centering
\caption{Survey-weight pseudo-posterior sensitivity}
\small
\begin{tabular}{lrrr}
\toprule
Comparison & Median OR & 2.5th percentile & 97.5th percentile \\
\midrule
Male vs female & 1.53 & 1.22 & 1.97 \\
Age 24--35 vs 0--5 months & 3.45 & 2.20 & 5.38 \\
Richest vs poorest & 0.30 & 0.16 & 0.61 \\
Unimproved water & 1.38 & 0.90 & 2.07 \\
Unimproved sanitation & 1.14 & 0.82 & 1.65 \\
Higher maternal education vs none & 0.31 & 0.15 & 0.61 \\
\bottomrule
\end{tabular}
\begin{minipage}{0.98\textwidth}\footnotesize \textit{Note.} These historical estimates are retained from Version 1. The separately audited local weighted fit passes the numerical screen but is not substituted for this table.\end{minipage}
\end{table}

\section{Missing maternal-education sensitivity}

A historical full-sample sensitivity model retained all 4,928 children by treating missing maternal education as a separate category. Its saved posterior was not among the nine audited files, so its joint diagnostics remain unverified. The main pattern was similar: unimproved water OR 1.35 (1.01--1.74), higher maternal education OR 0.38 (0.22--0.63), and a much less precise estimate for the missing-education category, OR 1.11 (0.80--1.53).


\begin{table}[H]
\centering
\caption{Full-sample missing-maternal-education sensitivity}
\small
\begin{tabular}{lrrr}
\toprule
Comparison & Median OR & 2.5th percentile & 97.5th percentile \\
\midrule
Male vs female & 1.44 & 1.21 & 1.70 \\
Age 24--35 vs 0--5 months & 2.90 & 2.11 & 4.10 \\
Richest vs poorest & 0.39 & 0.24 & 0.66 \\
Unimproved water & 1.35 & 1.01 & 1.74 \\
Unimproved sanitation & 1.15 & 0.90 & 1.48 \\
Missing maternal education vs none & 1.11 & 0.80 & 1.53 \\
Higher maternal education vs none & 0.38 & 0.22 & 0.63 \\
\bottomrule
\end{tabular}
\begin{minipage}{0.98\textwidth}\footnotesize \textit{Note.} Retained historical summaries. A saved posterior for this sensitivity was unavailable for the new full-parameter audit.\end{minipage}
\end{table}

\section{Age functional-form sensitivity}

The historical cubic B-spline summaries were similar to the categorical-age estimates, but the saved spline reproduction fit failed the strict joint screen. The following values remain exploratory:

\begin{table}[H]
\centering
\caption{Selected age-spline sensitivity results}
\begin{tabular}{lrrr}
\toprule
Comparison & Median OR & 2.5th percentile & 97.5th percentile \\
\midrule
Male vs female & 1.53 & 1.27 & 1.86 \\
Richest vs poorest & 0.35 & 0.19 & 0.61 \\
Unimproved water & 1.49 & 1.12 & 2.00 \\
Unimproved sanitation & 1.11 & 0.86 & 1.44 \\
Higher maternal education & 0.37 & 0.20 & 0.67 \\
\bottomrule
\end{tabular}
\end{table}

\section{Detailed WASH sensitivity}

\begin{table}[H]
\centering
\caption{Selected detailed-WASH sensitivity results}
\begin{tabular}{lrrr}
\toprule
Comparison & Median OR & 2.5th percentile & 97.5th percentile \\
\midrule
Surface water vs improved & 1.52 & 1.10 & 2.08 \\
Unprotected groundwater vs improved & 1.41 & 0.90 & 2.21 \\
Open defecation vs improved sanitation & 1.23 & 0.92 & 1.64 \\
Other unimproved vs improved sanitation & 0.93 & 0.68 & 1.31 \\
Higher maternal education vs none & 0.38 & 0.20 & 0.67 \\
\bottomrule
\end{tabular}
\end{table}

\section{Independent GEE robustness checks}

As an independent check, we fitted logistic generalized estimating equation (GEE) models with exchangeable working correlation at the household and community levels. The estimated working correlation was 0.154 for household clustering and 0.024 for community clustering. These values are not Bayesian ICCs; they are used only to check whether the stronger within-household dependence appears under a different modelling framework.

The GEE fixed-effect estimates were similar to the Bayesian sensitivity results. In the binary-WASH specification, unimproved water had OR 1.40 (1.12--1.74). Replacing categorical age with a cubic spline gave OR 1.39 (1.12--1.74). With detailed WASH coding, surface water had OR 1.43 (1.12--1.83), while the estimate for unprotected groundwater was less precise at OR 1.34 (0.93--1.94).

\begin{table}[H]
\centering
\caption{Selected independent GEE robustness checks}
\label{tab:gee_checks}
\begin{tabular}{lrrr}
\toprule
Specification / comparison & OR & Lower 95\% CI & Upper 95\% CI \\
\midrule
Binary WASH: unimproved water & 1.40 & 1.12 & 1.74 \\
Spline age: unimproved water & 1.39 & 1.12 & 1.74 \\
Detailed WASH: surface water & 1.43 & 1.12 & 1.83 \\
Detailed WASH: unprotected groundwater & 1.34 & 0.93 & 1.94 \\
\bottomrule
\end{tabular}
\end{table}

\section{Predictive checks and PSIS-LOO}

Posterior predictive checks reproduced the unweighted prevalence pattern in the 4,503-child complete-case Model 3 sample, both overall and across age groups. These PPC values are not the same quantity as the survey-weighted prevalence reported for the 4,928-child descriptive sample. For the strengthened Model 3 run, PSIS-LOO gave ELPD $-2011.47$, with maximum Pareto $k=0.69$ and no observations above 0.70. Harmonized Model 2 had five observations above Pareto $k=0.70$ and none above one. The ELPD difference favouring Model 3 was 0.92 with SE 3.38, so the retained comparison provides an inconclusive ranking and does not establish predictive equivalence. Joint-posterior limitations remain for the locked primary result and the audited harmonized comparator.


\begin{table}[H]
\centering
\caption{Retained exploratory observation-level PSIS-LOO comparison}
\small
\begin{tabular}{lrr}
\toprule
Quantity & Model 2 & Model 3 \\
\midrule
ELPD-LOO & -2012.39 & -2011.47 \\
Observations with Pareto $k>0.70$ & 5 & 0 \\
Observations with Pareto $k>1.00$ & 0 & 0 \\
\bottomrule
\end{tabular}
\begin{minipage}{0.98\textwidth}\footnotesize \textit{Note.} Model 3 minus Model 2 ELPD difference = 0.92 (SE 3.38). This conditional observation-level comparison remains inconclusive; it does not establish equivalence or performance in new households/communities, and joint-posterior diagnostic limitations remain.\end{minipage}
\end{table}

\section{Additional descriptive figures}

\begin{figure}[H]
\centering
\includesvg[width=0.80\textwidth]{stunting_by_sex}
\caption{Survey-weighted stunting prevalence by child sex.}
\end{figure}

\begin{figure}[H]
\centering
\includesvg[width=0.80\textwidth]{stunting_by_age}
\caption{Survey-weighted stunting prevalence by child age group.}
\end{figure}

\begin{figure}[H]
\centering
\includesvg[width=0.80\textwidth]{stunting_by_residence}
\caption{Survey-weighted stunting prevalence by place of residence.}
\end{figure}

\begin{figure}[H]
\centering
\includesvg[width=0.80\textwidth]{stunting_by_wealth}
\caption{Survey-weighted stunting prevalence by household wealth quintile.}
\end{figure}

\begin{figure}[H]
\centering
\includesvg[width=0.80\textwidth]{stunting_by_maternal_education}
\caption{Survey-weighted stunting prevalence by maternal education.}
\end{figure}

\begin{figure}[H]
\centering
\includesvg[width=0.80\textwidth]{stunting_by_water_source}
\caption{Survey-weighted stunting prevalence by drinking-water source.}
\end{figure}

\begin{figure}[H]
\centering
\includesvg[width=0.80\textwidth]{stunting_by_sanitation}
\caption{Survey-weighted stunting prevalence by sanitation category.}
\end{figure}

\section{Additional hierarchical-model figures}

\begin{figure}[H]
\centering
\includesvg[width=0.82\textwidth]{model_variance_comparison}
\caption{Household and community variance components across the principal model sequence.}
\end{figure}

\begin{figure}[H]
\centering
\includesvg[width=0.88\textwidth]{model3_ppc_age}
\caption{Posterior predictive check for unweighted observed prevalence in the 4,503-child complete-case Model 3 sample. These values are distinct from the survey-weighted prevalence reported for the full descriptive sample.}
\end{figure}

\section{Model-sequence variance summary}

\begin{table}[H]
\centering
\caption{Selected variance quantities across Models 1--3}
\resizebox{\textwidth}{!}{%
\begin{tabular}{lrrrrrr}
\toprule
Model & Community SD & Household SD & Community ICC & Household VPC & Community MOR & Household MOR \\
\midrule
Model 1 & 0.416 & 1.066 & 0.038 & 0.246 & 1.49 & 2.76 \\
Model 2 full sample & 0.390 & 1.100 & 0.033 & 0.260 & 1.45 & 2.86 \\
Model 2 harmonized (same 4,503-child sample as Model 3) & 0.391 & 1.222 & 0.031 & 0.301 & 1.45 & 3.21 \\
Model 3, locked eight-chain & 0.389 & 1.234 & 0.030 & 0.307 & 1.45 & 3.24 \\
\bottomrule
\end{tabular}%
}
\end{table}

\section{Selected final-model parameter diagnostics}

\begin{table}[H]
\centering
\caption{Selected parameter diagnostics from the locked eight-chain Model 3}
\resizebox{\textwidth}{!}{%
\begin{tabular}{lrrrrr}
\toprule
Parameter & Median & 2.5th percentile & 97.5th percentile & $\hat R$ & Bulk ESS \\
\midrule
Intercept & -2.598 & -3.276 & -1.955 & 1.003 & 1,830 \\
Community SD & 0.389 & 0.067 & 0.610 & 1.021 & 544 \\
Household SD & 1.234 & 0.940 & 1.535 & 1.018 & 524 \\
Male coefficient & 0.409 & 0.218 & 0.599 & 1.001 & 6,871 \\
Age 24--35 coefficient & 1.070 & 0.721 & 1.437 & 1.001 & 3,302 \\
Richest wealth coefficient & -1.020 & -1.622 & -0.449 & 1.000 & 3,670 \\
Unimproved-water coefficient & 0.404 & 0.138 & 0.687 & 1.001 & 5,444 \\
Unimproved-sanitation coefficient & 0.112 & -0.155 & 0.377 & 1.002 & 6,098 \\
Higher maternal education coefficient & -1.003 & -1.624 & -0.433 & 1.002 & 6,030 \\
\bottomrule
\end{tabular}%
}
\end{table}

\section{Full-posterior reproduction audit, 9 October 2026}

This post-preprint review read nine saved local fits without new sampling. It audited every stored parameter, including all household/community latent and scaled effects, and scientific quantities reconstructed from the SD draws: their difference, community ICC, household VPC, same-household ICC, and both MORs. Each full audit covered 7,832--8,359 scalar posterior elements; a supplemental calculation checked the six derived scientific quantities. The table reports the most conservative extrema across both checks. The screen required at least four chains, finite rank-normalized $\hat R<1.01$, bulk and tail ESS at least 400, zero divergences, minimum chain BFMI at least 0.3, and no stored maximum-tree-depth saturation flags.

\begin{table}[H]
\centering
\caption{Full-parameter diagnostics for separate four-chain reproduction fits}
\label{tab:full_posterior_audit}
\small
\begin{tabular}{lrrrrl}
\toprule
Fit & Max $\hat R$ & Min bulk ESS & Min tail ESS & Min BFMI & Screen \\
\midrule
Model 1 & 1.010885 & 485 & 533 & 0.496 & Fail \\
Model 2 full & 1.011620 & 651 & 1042 & 0.515 & Fail \\
Model 2 harmonized & 1.010417 & 568 & 534 & 0.509 & Fail \\
Spline age & 1.012751 & 447 & 612 & 0.532 & Fail \\
Detailed WASH & 1.015940 & 500 & 748 & 0.524 & Fail \\
Local Model 3 & 1.009845 & 456 & 514 & 0.482 & Pass* \\
Tight prior & 1.011931 & 456 & 620 & 0.516 & Fail \\
Wide prior & 1.011403 & 419 & 688 & 0.512 & Fail \\
Survey weighted & 1.008136 & 425 & 1065 & 0.762 & Pass* \\
\bottomrule
\end{tabular}

\begin{minipage}{0.98\textwidth}
\footnotesize
\textit{Note.} All nine fits had four chains, finite full-parameter diagnostics, zero divergences and zero tree-depth saturation flags. Seven fail the strict below-1.01 $\hat R$ rule. ESS and BFMI alone do not override a failure. *Pass is a numerical screen pending scientific review, not publication sign-off. Fits retained 1,000 draws per chain except survey weighting (1,500). These are separate reproduction fits; none is the locked eight-chain Model 3 used in the historical result tables.
\end{minipage}
\end{table}

Diagnostic calculation used Python 3.12 and ArviZ 0.22.0; archived file metadata reports PyMC 5.28.5 and ArviZ 0.23.4. The authorized analytic sample and all 41 descriptive rows were independently reproduced, with count agreement exactly and prevalence, SE and CI differences below $10^{-16}$. Stored observed outcomes matched the reconstructed full/complete sample for eight files; the weighted file did not store an observed-outcome array. Matching outcomes does not establish ordered child identity, exact original source code, seeds, or full fit provenance. A saved missing-maternal-education fit was not available in this set and remains unaudited. Trace/rank plots and estimand-specific Monte Carlo errors were not certified by this screen. Individual effects, coordinates and row-level arrays remain private.

The retained locked eight-chain SD diagnostics (community 1.021127 and household 1.017634) fail the same $\hat R$ screen. A passing local four-chain Model 3 does not certify that distinct locked fit. Neither the complete sensitivity sequence nor the historical conditional PSIS-LOO comparison is fully validated by the two numerical passes. No new WAIC/LOO ranking or model sampling was performed in this audit.

\end{document}
